\documentclass[10pt]{article}
\usepackage[margin=0.72in]{geometry}
\usepackage{amsmath,amssymb,mathtools}
\usepackage[T1]{fontenc}
\usepackage{lmodern}
\usepackage{microtype}
\usepackage{enumitem}
\setlength{\parindent}{0pt}
\setlength{\parskip}{4pt}
\setlist[itemize]{leftmargin=1.2em,itemsep=1pt,topsep=2pt}
\begin{document}
\small

\begin{center}
{\Large Phantom-aware shrinkage sampling: posterior construction and budget tradeoff}\\[2mm]
{\normalsize Concise note based on the design notes and reference implementation}
\end{center}

\section*{1. Posterior, global bootstrap, and the $\rho$ distribution}

Consider one classic boundary pair $a \to b$ with $b>a$. Among phantom samples drawn under constraints $c_i \le a$, define
\[
A = \#\{L>a\},\qquad B = \#\{L>b\},\qquad E = \#\{L=b\}.
\]
Here $A$ counts qualifiers from $\pi_a$, $B$ counts successes above the next boundary, and $E$ counts equality mass at the boundary. The design note uses a three-cell model
\[
p=(p_>,p_=,p_<),\qquad p \sim \mathrm{Dirichlet}(K,\varepsilon,1-\varepsilon),
\]
so that the shrinkage ratio $r = X(b)/X(a)=p_>$ has the classical marginal prior
\[
r \sim \mathrm{Beta}(K,1).
\]
Correlation among phantoms is handled by an efficiency parameter $\rho\in(0,1]$, interpreted as replacing the raw count $A$ by an effective count $A_{\mathrm{eff}}=\rho A$. Treating the observed proportions as fractional multinomial counts gives the posterior
\[
p\mid A,B,E,\rho \sim \mathrm{Dirichlet}\bigl(K+\rho B,\;\varepsilon+\rho E,\;1-\varepsilon+\rho(A-B-E)\bigr).
\]
Marginalizing out the equality and lower cells yields the shrinkage posterior used by the sampler,
\[
\boxed{\;r\mid A,B,\rho \sim \mathrm{Beta}(\alpha,\beta),\quad
\alpha=K+\rho B,\quad \beta=1+\rho(A-B).\;}
\]
When $A=0$ the code falls back to the classical $\mathrm{Beta}(K,1)$ update.

The global bootstrap is performed at the level of constrained-sampling calls rather than individual phantom points. Let cluster $i$ contain a correlated phantom set $\{L_{i,t}\}$ generated under constraint $c_i$. For bootstrap replicate $s$, sample clusters with replacement, keep each resampled cluster intact, and recompute all boundary counts $(A_g,B_g,E_g)$ from the same replicate. This preserves the within-cluster dependence and, crucially, the cross-boundary dependence induced by loose reuse: the same resampled cluster can contribute to many adjacent boundaries. Given the replicate counts, sample a single global $\rho_s$ and then sample all boundary shrinkages independently conditional on that replicate,
\[
r_{g,s}\sim \mathrm{Beta}\bigl(K_g+\rho_s B_g,\;1+\rho_s(A_g-B_g)\bigr),
\]
followed by the usual recursion
\[
X_{-1,s}=1,\qquad X_{g,s}=X_{g-1,s}r_{g,s},\qquad
Z_s=\sum_g \bigl(X_{g-1,s}-X_{g,s}\bigr)e^{\ell_g},
\]
with $\ell_g=\log L_{\text{block},g}$. This is the mechanism by which the sampler propagates both phantom reuse and $\rho$ uncertainty into $\log Z$.

The $\rho$ distribution comes from a candlestick calibration. For one boundary, define the observed proportion vector
\[
\hat p = (\hat p_>,\hat p_=)=\left(\frac{B}{A},\frac{E}{A}\right),
\]
and let $\mu$ and $\Sigma$ be the mean and covariance of the first two Dirichlet components under $\alpha=(K,\varepsilon,1-\varepsilon)$, with $\alpha_0=K+1$. In the iid case the covariance inflates like $\Sigma(1+\alpha_0/A)$; with correlated phantoms the design note replaces this by
\[
\kappa_g(\rho)=1+\frac{\alpha_{0,g}}{\rho A_g}.
\]
Using the Mahalanobis distance
\[
d_g^2=(\hat p_g-\mu_g)^\top \Sigma_g^{-1}(\hat p_g-\mu_g),
\]
and the approximation $d_g^2\sim \kappa_g(\rho)\chi_2^2$, the discrete grid posterior used in the code is
\[
\Pr(\rho=\rho_m\mid \text{bootstrap data}) \propto
\exp\!\Bigl[-\sum_{g:A_g>0}\Bigl\{\frac{2}{2}\log \kappa_g(\rho_m)+\frac{d_g^2}{2\kappa_g(\rho_m)}\Bigr\}\Bigr]\,\pi(\rho_m),
\]
with either $\pi(\rho)=1$ or the optional log prior $\pi(\rho)\propto 1/\rho$. Sampling $\rho$ from this grid rather than fixing it is what allows the evidence bootstrap to reflect uncertainty in phantom quality.

\section*{2. How to spend budget: more phantom samples or more live points?}

Write the chain length as $xD$ slice steps per replacement, of which $yD$ are burned in and discarded. Then the number of retained phantom points per replacement is
\[
p = xD-yD-1.
\]
If each slice costs about $n$ likelihood evaluations, then one replacement costs $nxD$ likelihood calls. Traversing the information-bearing region costs about $KH$ replacements, where $H$ is the usual nested-sampling information. Hence the total budget obeys
\[
N_{\mathrm{like}}\approx n\,xD\,K\,H.
\]

To summarize how much loose reuse actually helps, introduce a reuse efficiency $\eta\in[1/(K+1),1]$ defined so that the pooled effective phantom information behaves like $\rho\eta p$ extra shrinkage information per boundary. Then a useful evidence-level approximation is
\[
\boxed{\;\mathrm{Var}(\log Z)\approx \frac{H}{K_{\mathrm{ev}}},\qquad
K_{\mathrm{ev}}:=K\bigl(1+\rho\eta p\bigr).\;}
\]
The two important limits are:
\begin{itemize}
    \item fresh-only use of each replacement's phantoms: $\eta\approx 1/(K+1)$, so $K_{\mathrm{ev}}\approx K+\rho p$;
    \item strong loose reuse across many later boundaries: $\eta\approx 1$, so $K_{\mathrm{ev}}\approx K(1+\rho p)$.
\end{itemize}
Thus phantoms can improve evidence precision either additively or multiplicatively, depending on how much reuse survives in practice.

For a fair comparison at fixed likelihood budget, define
\[
C:=\frac{N_{\mathrm{like}}}{nDH},\qquad Kx=C,
\]
so increasing $x$ forces $K$ down. Using $p=xD-yD-1$, the evidence-effective live count becomes
\[
K_{\mathrm{ev}}(K)=K\Bigl[1+\rho\eta\Bigl(D\Bigl(\frac{C}{K}-y\Bigr)-1\Bigr)\Bigr]
= \rho\eta DC + \bigl[1-\rho\eta(Dy+1)\bigr]K.
\]
This produces a clean threshold:
\[
\boxed{\;\rho\eta(Dy+1)>1 \Longrightarrow \text{evidence prefers smaller $K$ and larger $x$};\;}
\]
\[
\boxed{\;\rho\eta(Dy+1)<1 \Longrightarrow \text{evidence prefers larger $K$ and smaller $x$}.\;}
\]
So, for evidence alone, there is no interior optimum in this simple model; the optimum sits on one side of the threshold.

Posterior quality pulls the other way. If only classic dead points are retained as posterior samples, then their reservoir size and effective sample size scale roughly like
\[
\mathrm{ESS}_{\mathrm{classic}} \propto K.
\]
Therefore a small-$K$, large-$x$ design can improve evidence precision while simultaneously reducing posterior ESS. The operational rule is:
\begin{itemize}
    \item choose $y$ as the smallest burn-in that gets the slice chain close enough to stationarity;
    \item ensure $K$ is large enough for mode finding, exploration robustness, and posterior ESS;
    \item once that minimum viable $K$ is set, spend the remaining budget on larger $x$ only if $\rho\eta(Dy+1)>1$.
\end{itemize}
In other words, extra burn-in beyond stationarity is always a tax, while extra post-burn-in phantom samples are worthwhile only when their effective information rate $\rho\eta$ is high enough to compensate for the loss in live points.

\end{document}
